\chapter{Modern Attention Techniques}
\label{ch:modern-attn}

The GPT-2 architecture we built in the previous chapter is a faithful reproduction of the 2019 paper.  Since then, several techniques have become standard practice in production language models.  Each one is a self-contained improvement: you can adopt them individually, and each has a clear paper trail.

In this chapter we replace learned positional embeddings with \textbf{Rotary Position Embeddings} (RoPE), swap LayerNorm for \textbf{RMSNorm}, add \textbf{Grouped-Query Attention} (GQA), enable \textbf{Flash Attention} via PyTorch's SDPA, switch the FFN activation from GELU to $\mathrm{relu}^2$, and add \textbf{QK normalization}.

All changes land in \code{microgpt/attention.py} (the new attention module) and an updated \code{microgpt/model.py}.


\section{RMSNorm}

\textbf{Paper:} Zhang \& Sennrich, ``Root Mean Square Layer Normalization'' (2019).

LayerNorm normalizes by subtracting the mean and dividing by the standard deviation:
\[
  \text{LayerNorm}(x) = \gamma \cdot \frac{x - \mu}{\sqrt{\sigma^2 + \epsilon}} + \beta
\]

RMSNorm drops the mean-centering step.  It normalizes by the root mean square only:
\[
  \text{RMSNorm}(x) = \gamma \cdot \frac{x}{\text{RMS}(x) + \epsilon}, \qquad
  \text{RMS}(x) = \sqrt{\frac{1}{d}\sum_{i=1}^{d} x_i^2}
\]

This is simpler (one fewer reduction op, no learnable bias $\beta$) and empirically works just as well.  LLaMA, Gemma, and most modern LLMs use RMSNorm.  PyTorch provides \code{F.rms\_norm} for an optimized implementation.

\filetag{microgpt/model.py}
\begin{lstlisting}
class RMSNorm(nn.Module):
    """Root Mean Square LayerNorm (Zhang & Sennrich, 2019).

    RMSNorm(x) = x / RMS(x) * gamma
    where RMS(x) = sqrt(mean(x^2) + eps)

    Compared to LayerNorm, RMSNorm skips the mean-centering step.
    This makes it simpler (fewer ops, no bias parameter) and empirically
    works just as well.  Used by LLaMA, Gemma, and most modern LLMs.
    """

    def __init__(self, n_embd: int, eps: float = 1e-6) -> None:
        super().__init__()
        self.weight = nn.Parameter(torch.ones(n_embd))
        self.eps = eps

    def forward(self, x: torch.Tensor) -> torch.Tensor:
        return F.rms_norm(x, self.weight.shape, self.weight, self.eps)
\end{lstlisting}


\section{Rotary Position Embeddings (RoPE)}

\textbf{Paper:} Su et al., ``RoFormer: Enhanced Transformer with Rotary Position Embedding'' (2021).

\subsection{The problem with learned positional embeddings}

BabyGPT and canonical GPT-2 add a learned position vector to each token:
\[
  x_t = \text{token\_emb}(t) + \text{pos\_emb}(t)
\]

This has two limitations:
\begin{enumerate}
  \item \textbf{Fixed maximum length.}  The position embedding table has \code{sequence\_len} rows.  The model cannot handle sequences longer than what it was trained on.
  \item \textbf{Absolute positions.}  Position 5 always gets the same vector, regardless of context.  The model has no built-in notion of ``3 tokens apart.''
\end{enumerate}

\subsection{The RoPE idea}

RoPE encodes position by \emph{rotating} query and key vectors in 2D subspaces.  Instead of adding a position vector, it multiplies by a rotation matrix.

Given a head dimension $d$, pair up the dimensions into $d/2$ pairs.  Each pair $(x_{2i}, x_{2i+1})$ is a 2D vector.  At position $m$, rotate this pair by angle $m \cdot \theta_i$:

\[
\begin{pmatrix} y_{2i} \\ y_{2i+1} \end{pmatrix}
=
\begin{pmatrix}
\cos(m\theta_i) & -\sin(m\theta_i) \\
\sin(m\theta_i) & \phantom{-}\cos(m\theta_i)
\end{pmatrix}
\begin{pmatrix} x_{2i} \\ x_{2i+1} \end{pmatrix}
\]

where $\theta_i = 1 / 10000^{2i/d}$.

The key insight: when computing $q_m^T k_n$ (the attention score between positions $m$ and $n$), the rotations combine as $\cos((m-n)\theta_i)$ --- the score depends only on the \emph{relative} distance $m - n$.  This gives the model relative position awareness with zero learned parameters.

\subsection{Implementation}

We precompute the $\cos$ and $\sin$ tables once and apply them in the attention forward pass.

\filetag{microgpt/attention.py}
\begin{lstlisting}
"""Modern attention components: RoPE, GQA, QK-norm, Flash Attention.

This module replaces the canonical GPT-2 attention (Chapter 2) with
techniques that have become standard since 2019.  Each is a self-contained
improvement with a clear paper trail.
"""

import math

import torch
import torch.nn as nn
import torch.nn.functional as F


# ---------------------------------------------------------------------------
# Rotary Position Embeddings (RoPE) --- Su et al. 2021
# ---------------------------------------------------------------------------
# Instead of adding a learned position vector to each token, RoPE encodes
# relative position by *rotating* query and key vectors in 2D subspaces.
#
# Given a head-dim d, we pair up the dimensions into d/2 pairs.
# Each pair (x_{2i}, x_{2i+1}) is treated as a 2D vector and rotated
# by an angle theta_i * position, where theta_i = 1 / 10000^{2i/d}.
#
# The rotation matrix for pair i at position m is:
#
#   [ cos(m*theta_i)  -sin(m*theta_i) ]   [ x_{2i}   ]
#   [ sin(m*theta_i)   cos(m*theta_i) ] * [ x_{2i+1} ]
#
# This means: attention between positions m and n depends only on the
# *difference* (m - n), giving the model relative position awareness
# without any learned parameters.


def precompute_rope_frequencies(
    head_dim: int,
    max_seq_len: int,
    theta: float = 10000.0,
    device: torch.device | None = None,
) -> tuple[torch.Tensor, torch.Tensor]:
    """Precompute cos/sin tables for RoPE.

    Returns (cos, sin) each of shape (1, max_seq_len, 1, head_dim/2).
    The singleton dims broadcast over batch and heads.
    """
    assert head_dim % 2 == 0
    # theta_i = 1 / 10000^{2i/d} for i = 0, 1, ..., d/2 - 1
    half = head_dim // 2
    freqs = 1.0 / (theta ** (torch.arange(0, head_dim, 2, device=device).float() / head_dim))
    # positions: 0, 1, ..., max_seq_len - 1
    positions = torch.arange(max_seq_len, device=device).float()
    # outer product: (max_seq_len, head_dim/2)
    angles = torch.outer(positions, freqs)
    # Shape: (1, max_seq_len, 1, d/2) --- broadcasts over (B, T, H, D/2)
    cos = torch.cos(angles).unsqueeze(0).unsqueeze(2)  # (1, max_seq_len, 1, d/2)
    sin = torch.sin(angles).unsqueeze(0).unsqueeze(2)
    return cos, sin
\end{lstlisting}

The application function splits each head into two halves and rotates:

\filetag{microgpt/attention.py}
\begin{lstlisting}
def apply_rotary_emb(
    x: torch.Tensor,
    cos: torch.Tensor,
    sin: torch.Tensor,
) -> torch.Tensor:
    """Apply RoPE rotation to x.

    x: (B, T, H, D)  where D = head_dim
    cos, sin: (1, T, 1, D/2)  --- precomputed tables
    """
    d = x.shape[-1] // 2
    x1, x2 = x[..., :d], x[..., d:]
    # Rotation: [x1, x2] -> [x1*cos + x2*sin, -x1*sin + x2*cos]
    # Wait --- the standard rotation is [x1*cos - x2*sin, x1*sin + x2*cos].
    # The nanochat convention (which we follow) swaps the sign pattern.
    # Both are valid --- they just rotate in opposite directions.
    y1 = x1 * cos + x2 * sin
    y2 = x1 * (-sin) + x2 * cos
    return torch.cat([y1, y2], dim=-1)
\end{lstlisting}

\tip{There are two common implementations of RoPE.  The ``split'' version (used here and by nanochat) puts the first $d/2$ dims in one group and the second $d/2$ in another.  The ``interleaved'' version (used by LLaMA) pairs up even and odd indices.  Both produce the same rotary behavior --- the dot product between rotated $q$ and $k$ depends on relative position regardless of which pairing you use.}


\section{Grouped-Query Attention (GQA)}

\textbf{Paper:} Ainslie et al., ``GQA: Training Generalized Multi-Query Transformers from Multi-Head Checkpoints'' (2023).

Standard multi-head attention (MHA) gives each head its own Q, K, and V projections.  With 12 heads, that means 12 independent sets of K and V matrices.

GQA reduces the number of K/V heads while keeping the full number of Q heads.  Groups of query heads share one K/V head:

\begin{center}
\begin{tabular}{lccc}
\toprule
Method & Q heads & KV heads & KV memory \\
\midrule
MHA & 12 & 12 & 1x \\
GQA (4 KV) & 12 & 4 & 0.33x \\
MQA & 12 & 1 & 0.08x \\
\bottomrule
\end{tabular}
\end{center}

The main benefit is \textbf{inference efficiency}: the KV cache (Chapter 7) shrinks proportionally to the number of KV heads.  With GQA-4 instead of MHA, the KV cache is 3x smaller, enabling longer sequences or larger batch sizes.

In the implementation, we project K and V to fewer dimensions, then expand them to match Q during the attention computation:

\filetag{microgpt/attention.py}
\begin{lstlisting}
class CausalSelfAttention(nn.Module):
    """Modern causal self-attention.

    Improvements over the canonical GPT-2 attention:
    - RoPE instead of learned positional embeddings
    - Grouped-Query Attention (GQA): n_kv_head <= n_head
    - QK normalization for training stability
    - PyTorch SDPA (Flash Attention / memory-efficient backend)
    """

    def __init__(
        self,
        n_embd: int,
        n_head: int,
        n_kv_head: int,
        bias: bool = False,
    ) -> None:
        super().__init__()
        assert n_embd % n_head == 0
        assert n_head % n_kv_head == 0

        self.n_head = n_head
        self.n_kv_head = n_kv_head
        self.head_dim = n_embd // n_head
        self.n_kv_groups = n_head // n_kv_head  # queries per KV head

        # Q projects to full n_head * head_dim; K and V project to n_kv_head * head_dim
        self.W_Q = nn.Linear(n_embd, n_head * self.head_dim, bias=bias)
        self.W_K = nn.Linear(n_embd, n_kv_head * self.head_dim, bias=bias)
        self.W_V = nn.Linear(n_embd, n_kv_head * self.head_dim, bias=bias)
        self.c_proj = nn.Linear(n_embd, n_embd, bias=bias)

    def forward(
        self,
        x: torch.Tensor,
        cos: torch.Tensor,
        sin: torch.Tensor,
    ) -> torch.Tensor:
        B, T, C = x.shape

        # Project to Q, K, V
        # Q: (B, T, n_head, head_dim)
        # K, V: (B, T, n_kv_head, head_dim)
        q = self.W_Q(x).view(B, T, self.n_head, self.head_dim)
        k = self.W_K(x).view(B, T, self.n_kv_head, self.head_dim)
        v = self.W_V(x).view(B, T, self.n_kv_head, self.head_dim)

        # Apply RoPE to queries and keys
        # cos/sin are (1, max_seq_len, 1, D/2); slice to (1, T, 1, D/2)
        q = apply_rotary_emb(q, cos[:, :T], sin[:, :T])
        k = apply_rotary_emb(k, cos[:, :T], sin[:, :T])

        # QK normalization: normalize Q and K after RoPE
        # Prevents attention logits from growing with depth.
        q = F.rms_norm(q, (self.head_dim,))
        k = F.rms_norm(k, (self.head_dim,))

        # Transpose to (B, H, T, D) for SDPA
        q = q.transpose(1, 2)  # (B, n_head, T, head_dim)
        k = k.transpose(1, 2)  # (B, n_kv_head, T, head_dim)
        v = v.transpose(1, 2)  # (B, n_kv_head, T, head_dim)

        # GQA: expand K and V heads to match Q heads
        # Each group of (n_head / n_kv_head) query heads shares one KV head.
        if self.n_kv_groups > 1:
            k = k.unsqueeze(2).expand(-1, -1, self.n_kv_groups, -1, -1)
            k = k.reshape(B, self.n_head, T, self.head_dim)
            v = v.unsqueeze(2).expand(-1, -1, self.n_kv_groups, -1, -1)
            v = v.reshape(B, self.n_head, T, self.head_dim)

        # Flash Attention via SDPA (automatically selects the best backend)
        out = F.scaled_dot_product_attention(q, k, v, is_causal=True)

        # Reassemble heads: (B, n_head, T, head_dim) -> (B, T, n_embd)
        out = out.transpose(1, 2).contiguous().view(B, T, C)
        return self.c_proj(out)
\end{lstlisting}


\section{QK Normalization}

\textbf{Papers:} Henry et al. (2020); Dehghani et al. (2023).

In deep transformers, attention logits ($q \cdot k / \sqrt{d}$) can grow large because the norms of $q$ and $k$ are unbounded.  Large logits push softmax toward a one-hot distribution, choking gradient flow.

QK normalization applies RMSNorm to queries and keys after projection (and after RoPE in our case).  This bounds the logit magnitude, stabilizing training without changing the attention mechanism.

The implementation is two lines in the attention forward method:
\begin{lstlisting}[language=Python]
q = F.rms_norm(q, (self.head_dim,))
k = F.rms_norm(k, (self.head_dim,))
\end{lstlisting}


\section{Flash Attention}

\textbf{Paper:} Dao et al., ``FlashAttention: Fast and Memory-Efficient Exact Attention with IO-Awareness'' (2022).

Standard attention materializes the full $T \times T$ attention matrix, requiring $O(T^2)$ memory.  Flash Attention computes the \emph{same} result using a tiled algorithm that never materializes the full matrix, reducing memory from $O(T^2)$ to $O(T)$ and running 2-4x faster through better GPU memory hierarchy utilization.

We do not implement Flash Attention ourselves.  PyTorch's \code{F.scaled\_dot\_product\_attention} (SDPA) automatically dispatches to the best available backend:
\begin{itemize}
  \item \textbf{FlashAttention v2} on Ampere+ GPUs (A100, RTX 3090, etc.)
  \item \textbf{Memory-efficient attention} (xformers) as fallback
  \item \textbf{Math-mode} (naive) as a last resort
\end{itemize}

The one-line call replaces our entire manual attention computation (scale, mask, softmax, dropout, matmul):
\begin{lstlisting}[language=Python]
out = F.scaled_dot_product_attention(q, k, v, is_causal=True)
\end{lstlisting}

\tip{If you are running on an H100 or newer GPU, you can install FlashAttention 3 (\code{flash-attn}) for even faster attention.  Nanochat uses FA3 when available and falls back to SDPA otherwise.  For this tutorial, SDPA is sufficient and works everywhere.}


\section{relu-squared Activation}

\textbf{Paper:} So et al., ``Primer: Searching for Efficient Transformers for Language Modeling'' (2021).

The GPT-2 FFN uses GELU:
\[
  \text{FFN}(x) = W_2 \cdot \text{GELU}(W_1 x)
\]

We replace GELU with $\text{relu}^2$:
\[
  \text{FFN}(x) = W_2 \cdot \text{relu}(W_1 x)^2
\]

Why?
\begin{itemize}
  \item \textbf{Sparsity.}  ReLU zeros out negative activations.  Squaring keeps them zero but amplifies positive activations, producing a naturally sparse hidden representation.
  \item \textbf{Simplicity.}  No special function --- just \code{F.relu(x).square()}.
  \item \textbf{Performance.}  Empirically competitive with GELU at GPT-2 scale.
\end{itemize}

\filetag{microgpt/model.py}
\begin{lstlisting}
class FeedForward(nn.Module):
    """Position-wise FFN with relu-squared activation (So et al., 2021).

    FFN(x) = W2 * relu(W1 * x)^2

    relu-squared is computationally simpler than GELU and produces sparse
    activations (many exact zeros plus a few strongly positive values).
    Empirically competitive with GELU at GPT-2 scale and above.
    """

    def __init__(self, config: GPTConfig) -> None:
        super().__init__()
        self.c_fc = nn.Linear(config.n_embd, 4 * config.n_embd, bias=config.bias)
        self.c_proj = nn.Linear(4 * config.n_embd, config.n_embd, bias=config.bias)

    def forward(self, x: torch.Tensor) -> torch.Tensor:
        x = self.c_fc(x)
        x = F.relu(x).square()  # relu^2
        x = self.c_proj(x)
        return x
\end{lstlisting}


\section{Putting It All Together}

Here is the updated model.  Compared to Chapter 2's canonical GPT-2:
\begin{itemize}
  \item \code{LayerNorm} $\to$ \code{RMSNorm}
  \item Learned \code{wpe} $\to$ RoPE (no position embedding table)
  \item Standard MHA $\to$ GQA with QK-norm and SDPA
  \item GELU $\to$ relu-squared
  \item Added \code{embed\_norm}: RMSNorm after token embedding (stabilizes early training)
\end{itemize}

\filetag{microgpt/model.py}
\begin{lstlisting}
class TransformerBlock(nn.Module):
    """Pre-norm transformer block with modern components.

    x = x + Attention(RMSNorm(x))
    x = x + FFN(RMSNorm(x))
    """

    def __init__(self, config: GPTConfig) -> None:
        super().__init__()
        self.ln_1 = RMSNorm(config.n_embd)
        self.attn = CausalSelfAttention(
            n_embd=config.n_embd,
            n_head=config.n_head,
            n_kv_head=config.n_kv_head,
            bias=config.bias,
        )
        self.ln_2 = RMSNorm(config.n_embd)
        self.ffn = FeedForward(config)

    def forward(
        self,
        x: torch.Tensor,
        cos: torch.Tensor,
        sin: torch.Tensor,
    ) -> torch.Tensor:
        x = x + self.attn(self.ln_1(x), cos, sin)
        x = x + self.ffn(self.ln_2(x))
        return x
\end{lstlisting}

\filetag{microgpt/model.py}
\begin{lstlisting}
class MicroGPT(nn.Module):
    """Decoder-only transformer language model (GPT-2 scale, modern arch).

    Key differences from the canonical GPT-2 (Chapter 2):
    - RMSNorm instead of LayerNorm
    - RoPE instead of learned positional embeddings (no wpe)
    - GQA support (n_kv_head can be < n_head)
    - QK normalization inside attention
    - relu-squared activation instead of GELU
    - Flash Attention via SDPA
    """

    def __init__(self, config: GPTConfig) -> None:
        super().__init__()
        self.config = config

        # Pad vocab to nearest multiple of 64 for tensor-core efficiency
        padded_vocab = ((config.vocab_size + 63) // 64) * 64

        self.transformer = nn.ModuleDict(dict(
            wte=nn.Embedding(padded_vocab, config.n_embd),
            blocks=nn.ModuleList(
                [TransformerBlock(config) for _ in range(config.n_layer)]
            ),
            ln_f=RMSNorm(config.n_embd),
        ))
        self.lm_head = nn.Linear(config.n_embd, padded_vocab, bias=False)

        # No positional embedding table --- RoPE handles positions
        # Precompute RoPE cos/sin tables
        head_dim = config.n_embd // config.n_head
        cos, sin = precompute_rope_frequencies(
            head_dim, config.sequence_len, device=None,
        )
        self.register_buffer("rope_cos", cos)
        self.register_buffer("rope_sin", sin)

        # RMSNorm after token embedding (stabilizes early training)
        self.embed_norm = RMSNorm(config.n_embd)

        # Initialize weights
        self.apply(self._init_weights)
        # Scaled init for residual projections
        for name, p in self.named_parameters():
            if name.endswith("c_proj.weight"):
                nn.init.normal_(p, mean=0.0, std=0.02 / math.sqrt(2 * config.n_layer))

    def _init_weights(self, module: nn.Module) -> None:
        if isinstance(module, nn.Linear):
            nn.init.normal_(module.weight, mean=0.0, std=0.02)
            if module.bias is not None:
                nn.init.zeros_(module.bias)
        elif isinstance(module, nn.Embedding):
            nn.init.normal_(module.weight, mean=0.0, std=0.02)

    def forward(
        self,
        tokens: torch.Tensor,
        targets: torch.Tensor | None = None,
    ) -> tuple[torch.Tensor, torch.Tensor | None]:
        B, T = tokens.shape
        assert T <= self.config.sequence_len

        # Token embeddings only --- no positional embeddings (RoPE handles this)
        x = self.transformer.wte(tokens)
        x = self.embed_norm(x)

        # Pass cos/sin through the blocks for RoPE
        cos = self.rope_cos.to(x.dtype)
        sin = self.rope_sin.to(x.dtype)

        for block in self.transformer.blocks:
            x = block(x, cos, sin)

        x = self.transformer.ln_f(x)
        logits = self.lm_head(x)
        # Crop logits to actual vocab_size (remove padding tokens)
        logits = logits[..., :self.config.vocab_size]

        loss = None
        if targets is not None:
            loss = F.cross_entropy(
                logits.reshape(-1, logits.size(-1)),
                targets.reshape(-1),
            )
        return logits, loss
\end{lstlisting}

The \code{generate()} method is unchanged from Chapter 2 --- Chapter 7 will upgrade it with a KV cache.

\checkpoint{Verify the modernized model:}

\shelltag
\begin{lstlisting}[language=bash]
uv run python -c "
from microgpt.config import GPTConfig
from microgpt.model import MicroGPT
import torch
# MHA (default): 12 query heads, 12 KV heads
model = MicroGPT(GPTConfig())
print(f'MHA params: {model.count_parameters():,}')
# GQA: 12 query heads, 4 KV heads
model_gqa = MicroGPT(GPTConfig(n_kv_head=4))
print(f'GQA params: {model_gqa.count_parameters():,}')
x = torch.randint(0, 32768, (2, 64))
logits, loss = model(x[:, :-1], x[:, 1:])
print(f'Forward OK: loss {loss.item():.4f}')
"
\end{lstlisting}

\outputtag
\begin{lstlisting}[language={}]
MHA params: 135,286,272
GQA params: 125,849,088
Forward OK: loss 10.59
\end{lstlisting}

GQA with 4 KV heads saves about 10M parameters (the K and V projections shrink from $768 \times 768$ to $768 \times 256$ per block).  The bigger win comes at inference when the KV cache is 3x smaller.

\tip{Nanochat includes several more experimental techniques beyond what we cover here: value embeddings (ResFormer-style), per-layer residual scaling lambdas, smear gates, and the Muon optimizer.  These are active areas of research with less established paper trails.  We teach the techniques that have become industry-standard (RoPE, RMSNorm, GQA, Flash Attention) and leave the experimental ones as pointers for further reading.}


\section{Summary}

We upgraded every major component of the GPT-2 architecture:

\begin{center}
\begin{tabular}{lll}
\toprule
Component & Canonical GPT-2 & MicroGPT (modern) \\
\midrule
Normalization & LayerNorm & RMSNorm \\
Position encoding & Learned embeddings & RoPE \\
Attention & Standard MHA & GQA + QK-norm + SDPA \\
FFN activation & GELU & relu-squared \\
\bottomrule
\end{tabular}
\end{center}

Each change is well-motivated, well-documented in the literature, and independently beneficial.  The model is now ready for large-scale training --- which requires a data pipeline (next chapter) and distributed computing infrastructure (Chapter 5).
